% =====================================================================
%  Lebenslauf Harshil Gorasiya (Deutsch)
%  Kompilieren mit pdflatex (zweimal). Legen Sie photo.png in denselben Ordner.
%  Datei fuer den Versand umbenennen: Lebenslauf_Harshil_Gorasiya.pdf
%  Layout bewusst ATS-freundlich: Standardschrift, Klartext, keine Icons,
%  keine Fortschrittsbalken, einspaltiger Textfluss, eindeutige Ueberschriften.
% =====================================================================
\documentclass[10pt,a4paper]{article}

\usepackage[a4paper,left=1.7cm,right=1.7cm,top=1.1cm,bottom=1.35cm]{geometry}
\usepackage[T1]{fontenc}
\usepackage[utf8]{inputenc}
\usepackage{textcomp}
\usepackage[scaled=0.96]{helvet}
\renewcommand{\familydefault}{\sfdefault}
\IfFileExists{ngerman.ldf}{\usepackage[ngerman]{babel}}{}
\usepackage{microtype}
\usepackage{xcolor}
\usepackage{graphicx}
\usepackage{enumitem}
\usepackage{titlesec}
\usepackage{array}
\usepackage{fancyhdr}
\usepackage{lastpage}
\makeatletter

\definecolor{navy}{HTML}{1B365D}
\definecolor{mid}{HTML}{5A6472}
\definecolor{tint}{HTML}{EEF2F7}

\usepackage[unicode,colorlinks=true,urlcolor=navy,linkcolor=navy,
  pdftitle={Lebenslauf Harshil Gorasiya},
  pdfauthor={Harshil Gorasiya},
  pdfsubject={Lebenslauf: KI/ML-Engineering, LLM, RAG, KI-Agenten},
  pdfkeywords={Werkstudent, Masterarbeit, KI, Machine Learning, LLM, RAG, Python, PyTorch}]{hyperref}

\input{glyphtounicode}
\pdfgentounicode=1

% Layout bleibt in jeder Umgebung identisch (keine Silbentrennung, Flattersatz)
\hyphenpenalty=10000
\exhyphenpenalty=10000
\emergencystretch=3em
\setlength{\parindent}{0pt}
\setlength{\parskip}{0pt}
\raggedright
\newlength{\leftcol}\setlength{\leftcol}{2.45cm}
\newcommand{\gap}{0.5em}
\setlist[itemize]{leftmargin=*,labelindent=\dimexpr\leftcol+\gap\relax,labelwidth=0.55em,labelsep=0.4em,align=left,itemsep=0.6pt,topsep=1.5pt,parsep=0pt,partopsep=0pt,label={\color{navy}\textbullet}}


% Fusszeile
\pagestyle{fancy}
\fancyhf{}
\renewcommand{\headrulewidth}{0pt}
\renewcommand{\footrulewidth}{0pt}
\fancyfoot[C]{\footnotesize\color{mid}Harshil Gorasiya \textperiodcentered{} Lebenslauf \textperiodcentered{} Seite \thepage{} von \pageref{LastPage}}

% Abschnittsueberschriften
\titleformat{\section}{\normalfont\bfseries\color{navy}\normalsize}{}{0pt}
  {\MakeUppercase}[\vspace{-0.45em}{\color{navy}\rule{\linewidth}{0.7pt}}]
\titlespacing*{\section}{0pt}{11pt}{5pt}

% Hilfsmakros
\newcommand{\sep}{\hspace{0.45em}\textbar\hspace{0.45em}}
\newcommand{\dotsep}{\hspace{0.3em}\textperiodcentered\hspace{0.3em}}
\newcommand{\entry}[3]{%
  \begingroup\raggedright\leftskip=\dimexpr\leftcol+\gap\relax\noindent
  \hspace*{-\dimexpr\leftcol+\gap\relax}%
  \vtop to 0pt{\hsize=\dimexpr\leftcol+\gap\relax\raggedright\small\color{mid}\noindent #1\par\vss}%
  #2\par #3\par\endgroup\vspace{5pt}}
\newcommand{\etitle}[2]{\textbf{#1}\ifx&#2&\else\textcolor{mid}{\sep}#2\fi}
\newcommand{\stack}[1]{{\small\itshape\color{mid}#1}\par\vspace{5pt}}
\newcommand{\kpi}[2]{\begin{minipage}[t]{\linewidth}\centering{\large\bfseries\color{navy}#1}\\[-1pt]{\scriptsize\color{mid}#2}\end{minipage}}

\renewcommand\normalsize{\@setfontsize\normalsize{10.4}{12.6}}
\makeatother
\begin{document}
\normalsize

\noindent
\begin{minipage}[t]{\dimexpr\textwidth-3.05cm\relax}
  {\fontsize{25}{28}\selectfont\bfseries\color{navy}Harshil Gorasiya}\\[4pt]
  {\bfseries LLM-Systeme: Retrieval, KI-Agenten und Forschung zu kleinen Modellen}\\[1pt]
  {\color{navy}\bfseries Python\dotsep PyTorch\dotsep LLM\dotsep RAG\dotsep Computer Vision}\\[5pt]
  {\small
  Heilbronn, Deutschland\sep +49 155 6351 7346\sep
  \href{mailto:harshil.gorasiya.0011@gmail.com}{harshil.gorasiya.0011@gmail.com}\\[1pt]
  LinkedIn: \href{https://www.linkedin.com/in/harshil-gorasiya-054a58319}{linkedin.com/in/harshil-gorasiya-054a58319}\\[1pt]
  GitHub: \href{https://github.com/N0t-Harshil}{github.com/N0t-Harshil}\sep
  Portfolio: \href{https://harshil-0011.github.io/My-Portfolio/}{harshil-0011.github.io/My-Portfolio}}
\end{minipage}\hfill
\begin{minipage}[t]{2.7cm}
  \raggedleft
  \IfFileExists{photo.png}{\setlength{\fboxsep}{0pt}\setlength{\fboxrule}{0.6pt}%
    \fcolorbox{navy}{white}{\includegraphics[width=2.7cm]{photo.png}}}{}
\end{minipage}

\vspace{7pt}
\noindent\colorbox{tint}{\parbox{\dimexpr\textwidth-2\fboxsep\relax}{\small
  \textbf{Gesucht:} Werkstudentenstelle (20 Std./Woche im Semester)\dotsep
  Vollzeit in den Semesterferien\dotsep Berufseinstieg ab 2027\\[1pt]
  \textbf{Verfügbar:} ab sofort\sep
  \textbf{Aufenthaltstitel:} gültige Aufenthaltserlaubnis mit Arbeitserlaubnis}}

\vspace{6pt}
\noindent
\begin{tabular*}{\textwidth}{@{\extracolsep{\fill}}*{5}{>{\centering\arraybackslash}p{0.185\textwidth}}}
  \kpi{+8 Pkt.}{Genauigkeit: 71 \% auf 79 \%} &
  \kpi{–20 \%}{Trainingszeit (V100-GPUs)} &
  \kpi{2,4\texttimes}{Durchsatz gegenüber YOLOv5s} &
  \kpi{28–32 FPS}{Echtzeit-Objekterkennung auf CPU} &
  \kpi{8 Anbieter}{LLMs mit Failover (Ardan-CLI)} \\
\end{tabular*}

\section*{Kurzprofil}
Masterstudent an der Hochschule Heilbronn. Ich baue LLM-Anwendungen: lokale RAG-Systeme mit Quellenangaben, KI-Agenten mit Tool-Nutzung und Computer-Vision-Modelle für den Einsatz auf CPU. Dazu zwei ML-Praktika, 10 Monate Webentwicklung mit KI-Funktionen und ein Forschungsprojekt zu Reasoning-Modulen für Sprachmodelle.

\section*{Studium}
\entry{seit 03/2025}
  {\etitle{M.Sc. Software Engineering and Management}{Hochschule Heilbronn}}
  {\small Englischsprachiger Studiengang\sep Notendurchschnitt bisher: 1,9\sep voraussichtlicher Abschluss: 2027}
\entry{2020 -- 2024}
  {\etitle{B.Tech. Information Technology}{A.\,D. Patel Institute of Technology, Gujarat}}
  {\small Indien\sep Note: 1,9 (umgerechnet auf deutsche Skala)}

\section*{Projekte}
\entry{KI-Agent}
  {\etitle{Ardan-CLI: KI-Coding-Agent für das Terminal}{}}
  {\stack{TypeScript, Node.js, MCP, LSP}
  \begin{itemize}
    \item Agent mit Streaming, Plan-Modus, Sub-Agenten und Checkpoints; \textbf{35 Werkzeuge}.
    \item \textbf{8 LLM-Anbieter} mit automatischem Failover; \textbf{0 Laufzeit-Abhängigkeiten}, MCP- und LSP-Client selbst implementiert.
    \item \textbf{8,1 ms} Overhead pro Agent-Schritt (ohne Modellzeit); automatisierte Tests auf Linux, macOS und Windows.
  \end{itemize}}
\entry{Recherche-Engine}
  {\etitle{Local Perplex: lokaler KI-Recherche-Assistent mit Quellenbelegen}{}}
  {\stack{Python, FastAPI, SSE, SQLite, Ollama, Docker}
  \begin{itemize}
    \item Durchsucht Web und eigene Dokumente (PDF, Word, Excel) und streamt Antworten mit Zitaten.
    \item \textbf{26 LLM-Anbieter}, 7 Suchbackends, auf Wunsch vollständig offline betreibbar.
    \item \textbf{17\texttimes{} bessere Trefferqualität} als meine erste Version (BEIR SciFact); Antwortseite in ca. \textbf{200 ms}; Wissensgraph mit \textbf{77\texttimes{}} schnelleren Schreibvorgängen.
  \end{itemize}}
\entry{RAG-System}
  {\etitle{ANVITHA: lokale Wissensdatenbank mit belegten Antworten}{}}
  {\stack{Python, FastAPI, Ollama, SQLite, BM25}
  \begin{itemize}
    \item Macht PDF-, Text- und Markdown-Dateien zu einer vernetzten Wissensdatenbank; jede Antwort nennt die Seite der Originalquelle.
    \item Jeder Satz der Antwort wird gegen die Quellen geprüft; nicht belegte Aussagen werden entfernt.
    \item Fragen außerhalb des Themas werden in unter \textbf{0,3 s} abgelehnt, ohne Modellaufruf.
    \item Notizen als editierbares Markdown mit vollständigem Änderungsprotokoll; läuft vollständig lokal über Ollama.
  \end{itemize}}
\newpage
\entry{Masterarbeit\\(laufend)}
  {\etitle{Ekdant: Reasoning-Modul für Sprachmodelle}{}}
  {\stack{Hochschule Heilbronn, betreut durch einen Professor\sep PyTorch, Hugging Face}
  \begin{itemize}
    \item Paralleles Debattenmodul (4 Pfade, Konsens) als zusätzliche Stufe auf SmolLM2-135M; \textbf{+19,6 Mio. Parameter}.
    \item Kontrollierter Vergleich mit gleich aufgebauter Baseline auf medizinischer Frage-Antwort-Aufgabe: bestes Modell in allen Metriken vorn, mit knappem Abstand.
    \item Literaturbasis mit 66 Fachartikeln.
  \end{itemize}}
\entry{Weitere Projekte}
  {\etitle{KI-Videohochskalierung}{Python, PyTorch}}
  {\small Windows-Desktop-Anwendung für Video-Super-Resolution mit Modellauswahl; Real-ESRGAN-Varianten mit PSNR und SSIM verglichen.}
\entry{\strut}
  {\etitle{MNIST-MLOps}{Python, PyTorch, Flask, SQLite}}
  {\small Inkrementelles Training, Verwaltung von Modellversionen und Web-Dashboard.}

\section*{Berufserfahrung}
\entry{05/2024 -- 02/2025\\(10 Monate)}
  {\etitle{AI \& Web Solutions Developer}{Aani Brothers Infotech, Surat, Indien}}
  {\begin{itemize}
    \item Entwicklung webbasierter Anwendungen: Geschäftsanforderungen in Benutzeroberflächen und Abläufe übersetzt, auch mit KI-Funktionen.
    \item Nutzerforschung, Wireframing und Prototyping; Designs für bessere Usability optimiert.
  \end{itemize}}
\entry{12/2023 -- 04/2024\\(5 Monate)}
  {\etitle{Deep Learning Intern}{PD Avenue, Surat, Indien}}
  {\begin{itemize}
    \item CNN mit Residual-Blöcken (TensorFlow/Keras) für Alters- und Geschlechtsklassifikation auf \textbf{22.000+ Bildern}; Genauigkeit von \textbf{71 \% auf 79 \%} gesteigert.
    \item Trainingszeit auf V100-GPUs um \textbf{20 \%} gesenkt (Learning-Rate-Scheduling, größere Batchgröße), ohne Genauigkeitsverlust.
    \item Data-Leakage erkannt und behoben.
  \end{itemize}}
\entry{05/2023 -- 06/2023\\(2 Monate)}
  {\etitle{Machine Learning Intern}{Dotcom IoT LLP, Surat, Indien}}
  {\begin{itemize}
    \item YOLOv8n-Objekterkennung auf CPU-Edge-Gerät (Intel Core i7): 720p-Video mit \textbf{28–32 FPS}; \textbf{2,4\texttimes{} höherer Durchsatz} als YOLOv5s.
    \item \textbf{3.000+ Bilder} für fünf Fahrzeugklassen annotiert.
    \item Falsch-Positive bei schwachem Licht durch IoU- und NMS-Anpassung um \textbf{18 \%} reduziert.
  \end{itemize}}

\section*{Kenntnisse}
\entry{Programmierung}
  {\small Python, TypeScript, Node.js, SQL, Bash}{}
\entry{Engineering}
  {\small Git, GitHub Actions (CI/CD), Docker, Linux, pytest, REST-APIs, FastAPI, Flask, Streaming (SSE), JSON-RPC}{}
\entry{ML \& Vision}
  {\small PyTorch, TensorFlow/Keras, Hugging Face Transformers, YOLOv8, CNNs mit Residual-Blöcken, Objekterkennung, Super-Resolution (Real-ESRGAN), OpenCV, Datenannotation, Hyperparameter-Tuning, Learning-Rate-Scheduling, Overfitting-Analyse}{}
\entry{LLM \& Agenten}
  {\small RAG-Systeme, Retrieval (BM25, Embeddings), Wissensgraphen, Quellenprüfung für Antworten, Prompt Engineering, Fine-Tuning, kleine Sprachmodelle (SmolLM2), KI-Agenten mit Tool-Nutzung (ReAct), Model Context Protocol (MCP), Anthropic API, OpenAI-kompatible APIs, Ollama}{}
\entry{Evaluation \& Forschung}
  {\small A/B-Experimentdesign, Modellbewertung (Genauigkeit, Latenz, Durchsatz, PSNR/SSIM, Perplexität, ROUGE-L), Literaturrecherche}{}
\entry{Datenbanken \& UX}
  {\small SQLite, PostgreSQL, MongoDB, NumPy\sep Nutzerforschung, Wireframing, Prototyping, Usability}{}

\section*{Zertifikate}
\entry{2026}{\small Claude with the Anthropic API\sep Model Context Protocol: Advanced Topics (beide Anthropic)}{}
\entry{2025}{\small Machine Learning with Python (IBM)\sep Fundamentals of Data Science (Databricks)}{}

\section*{Sprachen}
\entry{\strut}{\small \textbf{Englisch:} fließend\sep \textbf{Deutsch:} A2–B1\sep \textbf{Hindi:} fließend\sep \textbf{Gujarati:} Muttersprache}{}

\vspace{2pt}
{\small\color{mid}Heilbronn, 05.11.2026}

\end{document}
